\FloatBarrier
\subsection{Corrected multi-family exact-rational stress test}
\label{sec:p4-multifamily}
We further varied the spectral fixture on a fixed interval grid: three deterministic strictly
positive spectral families (linear, quadratic, and alternating),
five dimensions $p=2,3,6,9,12$, the same twelve distinct-up-to-scale rational
contrast patterns per spectrum, and both eight and sixteen equal
rational subintervals. This produces 4,320 matched interval cases. An initial
contrast generator made its twelve nominal patterns scalar multiples;
we detected this defect and replaced it before this corrected run.
This correction was made after examining the first sweep, so this
expanded comparison is exploratory rather than an independently
preregistered test.
The numerator construction, coefficients, interval endpoints and
both methods use exact rational arithmetic. These are synthetic
classifier spectra, not an end-to-end certificate of the numerical
EEG preprocessing pipeline.

\begin{table}[ht]\centering\small
\caption{Exact rational whole-interval strict-sign certifications
on the corrected 4,320-case exploratory sweep. Opposite-sign rational
point witnesses invalidate a whole-interval constant strict sign for
88 unresolved cases. The extra Bernstein signs include thirteen
quadratic-spectrum dimension-twelve cases; other gains are spread
across the remaining tested spectra and dimensions.}
\label{tab:p4-multifamily}
\begin{tabular}{lr}\toprule
Method or diagnostic & Interval count\\\midrule
Midpoint-centered Horner certified & 4,212\\
Bernstein certified & 4,232\\
Bernstein only versus centered Horner & 20\\
Centered Horner only versus Bernstein & 0\\
Both unresolved with opposite-sign rational witnesses & 88\\
Both unresolved without such a witness & 0\\
\bottomrule\end{tabular}
\end{table}

The Bernstein fallback raises the
strict-sign count from 4,212 to 4,232 on this corrected
synthetic suite, a 0.463-percentage-point change in
coverage (20/4,320). It does not repair the 88 genuinely
sign-changing intervals and does not prove an asymptotic
speed advantage. In the earlier quadratic-spectrum 960-case
fixture, the advantage remains zero. This result narrows the
claim to a small reproducible advantage on selected synthetic
intervals under these corrected patterns; tests of rounding-safe floating implementations,
root isolation, runtime and more diverse polynomials remain
necessary before a numerical-algorithm journal submission.
